% Paper 2, §9 — DPO and offline preference optimization
% Anchors: kb/notes/post-training/dpo-and-offline.md
%          kb/excerpts/{rafailov2023-dpo,meng2024-simpo,ouyang2022-instructgpt}.md
% Cross-paper: cited from §10 (RLAIF feeding DPO) and §11 (RLVR as
%              orthogonal axis).

\section{DPO and offline preference optimization}
\label{sec:dpo}

The closed-form RL optimum \cref{eq:rlhf-optimum} that closes
\cref{sec:rlhf} is also the load-bearing equation of this one. Direct
Preference Optimization \citep{rafailov2023-dpo} reads
\cref{eq:rlhf-optimum} not as the \emph{target} of an RL loop but as the
\emph{definition} of a reward that can be inverted: solving for $r$ in
terms of $\pi^*$ and $\pi^{\mathrm{ref}}$ gives the reward as a
log-ratio of policies, and substituting that log-ratio into the
Bradley--Terry preference probability \cref{eq:rlhf-rm} cancels the
intractable partition function and yields a closed-form classification
loss on the same preference data the \glsterm{reward-model}{reward model} of
\cref{sec:rlhf-rm} was trained on \citep[\S 4]{rafailov2023-dpo}.\footnote{KB
excerpt: \texttt{kb/excerpts/rafailov2023-dpo.md\#sec-4}.} The thesis
of this section is that this single algebraic move converts post-
training from a PPO-class engineering problem (four model copies, three-
phase rollout/score/optimize, on-policy sampling) into a
fine-tuning-class engineering problem (two forward passes, one backward,
no sampler, no \glsterm{value}{value} head): \emph{same fixed point on the same data,
dramatically cheaper per step.} The variant family that has grown
around \glsterm{dpo}{DPO} since 2023 (\glsterm{ipo}{IPO}, \glsterm{kto}{KTO}, \glsterm{orpo}{ORPO}, \glsterm{simpo}{SimPO}, CPO) inherits this
structure and varies one of three knobs at a time --- the loss shape on
the implicit-reward gap, the reference dependence, and the length
normalization --- producing a small organized lattice of offline
preference methods.

\subsection{Reparameterization: from RL optimum to classification loss}
\label{sec:dpo-derivation}

\cref{sec:rlhf-ppo}'s Eq.~\cref{eq:rlhf-optimum} states that the
KL-constrained reward maximization
$\max_{\pi_\theta} \E[r(x,y)] - \beta\,\KL(\pi_\theta \| \pi^{\mathrm{ref}})$
admits the analytic optimum
\begin{equation}
\pi^*(y \mid x)
  = \frac{1}{Z(x)} \, \pi^{\mathrm{ref}}(y \mid x)
    \exp\!\left(\frac{1}{\beta}\, r(x, y)\right),
\label{eq:dpo-rl-optimum}
\end{equation}
with $Z(x) = \sum_y \pi^{\mathrm{ref}}(y \mid x) \exp(r(x,y)/\beta)$ ---
a re-weighting of the reference distribution by an exponential of the
reward \citep[\S 4, Eq.~4]{rafailov2023-dpo}.\footnote{KB excerpt:
\texttt{kb/excerpts/rafailov2023-dpo.md\#sec-4}.} \citet{rafailov2023-dpo}'s
move is to invert this expression for $r$. Taking logarithms of
\cref{eq:dpo-rl-optimum} and rearranging,
\begin{equation}
r(x, y)
  = \beta \log \frac{\pi^*(y \mid x)}{\pi^{\mathrm{ref}}(y \mid x)}
    \;+\; \beta \log Z(x),
\label{eq:dpo-reward-from-policy}
\end{equation}
verbatim from \citet[\S 4, Eq.~5]{rafailov2023-dpo}. The reward is now
written as a function of the optimal policy and the reference, plus a
prompt-dependent additive constant $\beta \log Z(x)$. The constant is
intractable in general --- $Z(x)$ is a sum over all possible
completions $y$ --- but the next step removes it. The Bradley--Terry
preference assumption from \cref{sec:rlhf-rm} states
$p(y_w \succ y_l \mid x) = \sigma(r(x, y_w) - r(x, y_l))$, a difference
of rewards on the same prompt. Substituting
\cref{eq:dpo-reward-from-policy} into the Bradley--Terry probability,
\begin{equation}
p^*(y_w \succ y_l \mid x)
  = \sigma\!\left(
       \beta \log \frac{\pi^*(y_w \mid x)}{\pi^{\mathrm{ref}}(y_w \mid x)}
       - \beta \log \frac{\pi^*(y_l \mid x)}{\pi^{\mathrm{ref}}(y_l \mid x)}
     \right),
\label{eq:dpo-bt-pref}
\end{equation}
since the $\beta \log Z(x)$ terms appear with opposite signs on $y_w$
and $y_l$ and cancel exactly \citep[\S 4, Eq.~6]{rafailov2023-dpo}.
The preference probability now depends only on policy log-ratios; the
intractable normalizer is gone.

Maximum-likelihood on a preference dataset
$\mathcal{D} = \{(x_i, y_w^i, y_l^i)\}_{i=1}^N$ then yields the \glsterm{dpo}{DPO}
loss \citep[\S 4, Eq.~7]{rafailov2023-dpo}:\footnote{KB excerpt:
\texttt{kb/excerpts/rafailov2023-dpo.md\#sec-4-loss}.}
\begin{equation}
\mathcal{L}_{\mathrm{DPO}}(\pi_\theta;\, \pi^{\mathrm{ref}})
  = -\, \E_{(x, y_w, y_l) \sim \mathcal{D}}
      \!\left[
        \log \sigma\!\left(
          \beta \log \frac{\pi_\theta(y_w \mid x)}{\pi^{\mathrm{ref}}(y_w \mid x)}
          - \beta \log \frac{\pi_\theta(y_l \mid x)}{\pi^{\mathrm{ref}}(y_l \mid x)}
        \right)
      \right].
\label{eq:dpo-loss}
\end{equation}
The \glsterm{implicit-reward}{implicit reward} associated with the policy under training is
\begin{equation}
\hat{r}_\theta(x, y)
  \;=\; \beta \log \frac{\pi_\theta(y \mid x)}{\pi^{\mathrm{ref}}(y \mid x)},
\label{eq:dpo-implicit-reward}
\end{equation}
so \cref{eq:dpo-loss} is exactly the Bradley--Terry log-likelihood
\cref{eq:rlhf-rm} of \cref{sec:rlhf}, evaluated on the \glsterm{implicit-reward}{implicit reward}
\cref{eq:dpo-implicit-reward} instead of a separately-trained
$r_\phi$ \citep[\S 4]{rafailov2023-dpo}.

\paragraph{Same fixed point as PPO-RLHF.} The crucial structural fact:
if $\pi_\theta$ approaches the maximum-likelihood solution of
\cref{eq:dpo-loss}, then by construction the \glsterm{implicit-reward}{implicit reward}
\cref{eq:dpo-implicit-reward} approaches the underlying preference-data
reward $r$, and $\pi_\theta$ approaches the closed-form RL optimum
$\pi^*$ of \cref{eq:dpo-rl-optimum} on that reward. \emph{\glsterm{dpo}{DPO} and
PPO-\glsterm{rlhf}{RLHF} share the same fixed point on the same preference data}
\citep[\S 4]{rafailov2023-dpo} --- they differ only in the optimization
trajectory, not in the optimum. PPO walks toward $\pi^*$ along on-policy
rollouts scored by a learned $r_\phi$; \glsterm{dpo}{DPO} walks toward $\pi^*$ along
gradient steps of a classification loss on a fixed preference dataset.

\paragraph{The DPO gradient.} Differentiating
\cref{eq:dpo-loss}\footnote{KB excerpt:
\texttt{kb/excerpts/rafailov2023-dpo.md\#sec-4-gradient}.} gives
\citep[\S 4]{rafailov2023-dpo}
\begin{equation}
\nabla_\theta \mathcal{L}_{\mathrm{DPO}}
  = -\, \beta\, \E\!\left[\,
       \sigma\!\big(\hat{r}_\theta(x, y_l) - \hat{r}_\theta(x, y_w)\big)
       \;\big(\nabla_\theta \log \pi_\theta(y_w \mid x)
              - \nabla_\theta \log \pi_\theta(y_l \mid x)\big)
     \right],
\label{eq:dpo-gradient}
\end{equation}
the \emph{preference-margin gradient} that defines all \glsterm{dpo}{DPO}-class
methods. Two structural pieces are worth naming. First, the gradient is
proportional to a \emph{relative} score change: the policy is asked to
push up $\log \pi_\theta(y_w \mid x)$ and push down
$\log \pi_\theta(y_l \mid x)$ \emph{by the same amount}. Second, the
sigmoid weight $\sigma(\hat r_\theta(y_l) - \hat r_\theta(y_w))$ is
large precisely when the preference is currently violated by the
\glsterm{implicit-reward}{implicit reward} (the policy assigns lower \glsterm{implicit-reward}{implicit reward} to the
preferred response), and small when the policy already agrees with the
preference. The update is self-modulating: aggressive on
miscalibrated pairs, gentle on already-correct ones.

\begin{intuition}
The \glsterm{dpo}{DPO} derivation is an \emph{inversion} trick. The RL optimum
\cref{eq:dpo-rl-optimum} maps a reward $r$ to a policy $\pi^*$;
\cref{eq:dpo-reward-from-policy} reads the same equation backward and
maps a policy back to a reward. Once you have a policy-to-reward map,
the obvious move is to express the Bradley--Terry preference
probability \cref{eq:rlhf-rm} in terms of the policy directly, and the
intractable $\log Z(x)$ vanishes because it appears symmetrically on
$y_w$ and $y_l$. What looked like an RL problem (KL-constrained reward
maximization with on-policy rollouts) becomes a classification problem
(maximum-likelihood on policy log-ratios). Returning to canonical form:
\cref{eq:dpo-loss} is literally \cref{eq:rlhf-rm}, the
Bradley--Terry RM training loss, with the parametric $r_\phi$
substituted by the implicit $\hat r_\theta = \beta \log \pi_\theta /
\pi^{\mathrm{ref}}$ of \cref{eq:dpo-implicit-reward}. The ``language
model is secretly a \glsterm{reward-model}{reward model}'' phrasing
\citep{rafailov2023-dpo} is the algebraic identity, not a metaphor.
\end{intuition}

\subsection{The DPO training protocol}
\label{sec:dpo-protocol}

The protocol is supervised classification on log-ratios. Starting from
the SFT-stage model $\pi^{\mathrm{SFT}}$ of \cref{sec:sft}, set
$\pi^{\mathrm{ref}} \leftarrow \pi^{\mathrm{SFT}}$ and
$\pi_\theta \leftarrow \pi^{\mathrm{SFT}}$, then for each batch of
preference triples $(x, y_w, y_l)$:
\begin{enumerate}
\item Compute $\log \pi_\theta(y_w \mid x)$ and
$\log \pi_\theta(y_l \mid x)$ via a forward pass on the policy under
training.
\item Compute $\log \pi^{\mathrm{ref}}(y_w \mid x)$ and
$\log \pi^{\mathrm{ref}}(y_l \mid x)$ via a forward pass on the frozen
reference (often pre-cached, since $\pi^{\mathrm{ref}}$ does not
change).
\item Evaluate \cref{eq:dpo-loss} and back-propagate through $\theta$.
\end{enumerate}
There is no sampling, no \glsterm{reward-model}{reward model}, no \glsterm{value}{value} model, no reinforcement
learning loop. The total compute per step is $\sim\!2$ forward passes
$+ 1$ backward pass on $\pi_\theta$ (both win and loss completions
share the prompt prefix in practice, so the forward cost is closer to
$1.5\times$ a single completion's), versus PPO's four-model-copy plus
sampling-plus-scoring-plus-optimization three-phase loop
\citep[\S 3]{ouyang2022-instructgpt}. \cref{tab:dpo-vs-ppo-cost}
summarizes the per-step infrastructure delta that drives the
``\glsterm{dpo}{DPO}-replaces-PPO'' move at frontier scale.

\begin{table}[h]
\centering
\small
\begin{tabular}{lll}
\toprule
Cost axis & PPO-\glsterm{rlhf}{RLHF} & \glsterm{dpo}{DPO} \\
\midrule
Co-resident model copies & 4 ($\pi_\theta$, $\pi^{\mathrm{SFT}}$, $r_\phi$, $V_\phi$) & 2 ($\pi_\theta$, $\pi^{\mathrm{ref}}$) \\
Sampling phase & yes (rollouts at temperature $T$) & no \\
Scoring phase & yes ($r_\phi$ forward per rollout) & no \\
Forward passes per step & $\geq 4$ (rollout + $r_\phi$ + $V_\phi$ + ref) & 2 \\
Backward passes per step & 1--$K$ (PPO inner-loop $K \in \{1,4\}$) & 1 \\
Hyperparameters that matter & $\beta, \epsilon, \lambda_{\mathrm{GAE}}, \gamma_{\mathrm{ptx}}, \mathrm{LR}_{\pi}, \mathrm{LR}_V, T$ & $\beta, \mathrm{LR}$ \\
\bottomrule
\end{tabular}
\caption{Per-step cost comparison of PPO-RLHF (\cref{sec:rlhf-ppo}) vs.\
DPO (\cref{sec:dpo-protocol}). The ``four model copies'' line is the
load-bearing infrastructure fact behind the empirical move at frontier
scale \citep{meta-llama3, tulu3}.}
\label{tab:dpo-vs-ppo-cost}
\end{table}

\paragraph{Hyperparameters.} The canonical open-source recipe
\citep{tulu3} runs $\beta = 0.1$ and learning rate
$5 \times 10^{-7}$ --- about ten times smaller than the SFT learning
rate of \cref{sec:sft}. Llama-3 \citep{meta-llama3} uses the same
$\beta = 0.1$ at frontier scale. Batch sizes are typically
$32$--$256$ effective; epochs $1$--$3$. The low learning rate is not
optional: although \glsterm{dpo}{DPO} does not have PPO's \glsterm{value}{value}-network instability,
the policy can still drift far enough from $\pi^{\mathrm{ref}}$ to
collapse onto a degenerate mode that satisfies the preference margin on
training pairs while losing fluency on held-out prompts. The
$\mathrm{KL}(\pi_\theta \| \pi^{\mathrm{ref}})$ that emerges
implicitly --- through the $\hat r_\theta = \beta \log \pi_\theta /
\pi^{\mathrm{ref}}$ structure --- is the only thing keeping the policy
on the SFT support, and a higher LR overshoots that constraint.

\paragraph{On-policy preferences with off-policy optimization.} \glsterm{dpo}{DPO}'s
loss is off-policy by default: it reads a fixed dataset of preference
triples and never samples from $\pi_\theta$. This is structurally
identical to a one-shot SFT pass, but the empirical best practice at
frontier scale (Llama-3 \citep{meta-llama3}, T\"ulu-3 \citep{tulu3}) is
to \emph{iterate}: run \glsterm{dpo}{DPO} to convergence, sample new responses from
the resulting policy, label them (human or \glsterm{llm}{LLM}-judge) into new
preference triples, and run \glsterm{dpo}{DPO} again --- often with
$\pi^{\mathrm{ref}}$ kept at the original $\pi^{\mathrm{SFT}}$ across
rounds, so the \glsterm{implicit-reward}{implicit reward} stays anchored to the SFT support. The
result is \emph{on-policy preference collection with off-policy
optimization}: the gradient signal comes from preferences on
$\pi_\theta$'s own samples, recovering some of the on-policy benefit
of PPO without paying the per-step rollout cost.

\subsection{Variant family: one organized axis at a time}
\label{sec:dpo-variants}

The \glsterm{dpo}{DPO} variant landscape since 2023 has not produced a clear winner;
it has produced a small lattice of methods that vary one knob at a time
along three axes (loss shape on the implicit-reward gap; reference
dependence; length normalization). \cref{tab:dpo-variants} is the
load-bearing organizing table.

\begin{table}[h]
\centering
\small
\begin{tabular}{lcccc}
\toprule
Method & Reference & Reward form & Length-normalize & Loss shape \\
\midrule
PPO-\glsterm{rlhf}{RLHF} & $\pi^{\mathrm{SFT}}$ & learned $r_\phi$ & no & on-policy clipped surrogate \\
\glsterm{dpo}{DPO} & $\pi^{\mathrm{ref}}$ frozen & $\beta \log \pi_\theta / \pi^{\mathrm{ref}}$ & no & log-sigmoid \\
\glsterm{ipo}{IPO} & $\pi^{\mathrm{ref}}$ frozen & same as \glsterm{dpo}{DPO} & no & squared-gap \\
\glsterm{kto}{KTO} & $\pi^{\mathrm{ref}}$ frozen & prospect-theory utility & no & asymmetric per-rating \\
\glsterm{orpo}{ORPO} & none & SFT loss + log-odds-ratio & no & combined SFT + OR \\
\glsterm{simpo}{SimPO} & none & $(\beta/|y|) \log \pi_\theta(y \mid x)$ & yes & log-sigmoid + margin $\gamma$ \\
CPO & none & policy log-prob + SFT regularizer & no & log-sigmoid + SFT \\
\bottomrule
\end{tabular}
\caption{Offline preference variant family, organized by reference
dependence (anchored vs.\ reference-free) and the implicit-reward form.
PPO-RLHF row from \cref{sec:rlhf}; DPO row from
\citep[\S 4]{rafailov2023-dpo}; SimPO row from
\citep[\S 3]{meng2024-simpo}.\protect\footnotemark{} IPO, KTO, ORPO,
CPO entries summarize the canonical statement of each variant; see
\cref{sec:dpo-variants} for citations.}
\label{tab:dpo-variants}
\end{table}
\footnotetext{KB excerpts: \texttt{kb/excerpts/rafailov2023-dpo.md\#sec-4-loss},
\texttt{kb/excerpts/meng2024-simpo.md\#sec-3}.}

\paragraph{IPO --- squared-gap loss.} The \glsterm{dpo}{DPO} log-sigmoid loss has a
known overconfidence pathology: when one preference is clearly
satisfied, the \glsterm{implicit-reward}{implicit reward} gap $\hat r_\theta(y_w) - \hat
r_\theta(y_l)$ can grow without bound while the loss falls toward zero,
driving the policy to overfit a few extreme pairs and undertrain on the
bulk. \glsterm{ipo}{IPO} replaces the log-sigmoid with a squared loss on the gap,
\begin{equation}
\mathcal{L}_{\mathrm{IPO}}(\theta)
  = \E_{(x, y_w, y_l) \sim \mathcal{D}}
    \!\left[\,\Big(\hat{r}_\theta(x, y_w) - \hat{r}_\theta(x, y_l)
                   - \tfrac{1}{2\tau}\Big)^{\!2}\right],
\label{eq:ipo}
\end{equation}
where $\tau$ controls the target margin
\citep[Azar et al.\ 2023; Eq.~17 in the paper's IPO derivation]{rafailov2023-dpo}.\footnote{\glsterm{ipo}{IPO}'s
canonical form as reproduced in \glsterm{dpo}{DPO} survey arXiv:2503.11701 \\glsterm{s4}{S4}.1.}
\deepencite{azar2023-ipo §4 — bibkey not yet in references.bib}
The squared loss is bounded above and saturates symmetrically, removing
the runaway-log-ratio failure mode.

\paragraph{KTO --- single-rating data, prospect theory.} \glsterm{kto}{KTO} operates
on \emph{single-rating} data $(x, y, \mathrm{label})$ with
$\mathrm{label} \in \{\textsf{good}, \textsf{bad}\}$, instead of the
pairwise $(x, y_w, y_l)$ \glsterm{dpo}{DPO} requires. The loss is asymmetric, weighted
by Kahneman--Tversky prospect-theory loss-aversion --- a unit of
``bad'' reward weighs more than a unit of ``good'' reward
\citep[Ethayarajh et al.\ 2024]{rafailov2023-dpo}.
\deepencite{ethayarajh2024-kto §3 — bibkey not yet in references.bib}
\glsterm{kto}{KTO} is the right method when the labeling pipeline can produce thumbs-
up / thumbs-down judgments but not pairwise rankings (e.g., binary
production-feedback streams).

\paragraph{ORPO --- combined SFT and preference, no reference.} \glsterm{orpo}{ORPO}
collapses SFT and preference learning into a single objective and
removes the reference model entirely:
\begin{equation}
\mathcal{L}_{\mathrm{ORPO}}(\theta)
  = \mathcal{L}_{\mathrm{SFT}}(\pi_\theta;\, y_w)
    \;+\; \lambda \, \mathcal{L}_{\mathrm{OR}}(\pi_\theta;\, y_w, y_l),
\label{eq:orpo}
\end{equation}
where $\mathcal{L}_{\mathrm{SFT}}$ is the cross-entropy on the chosen
response (Eq.~\cref{eq:rlhf-sft}) and $\mathcal{L}_{\mathrm{OR}}$ is a
log-odds-ratio penalty between $y_w$ and $y_l$
\citep[Hong et al.\ 2024]{rafailov2023-dpo}.
\deepencite{hong2024-orpo §3 — bibkey not yet in references.bib}
The SFT term takes over the role of the KL anchor (without a separate
reference forward), and the pipeline reduces from two stages
(SFT $\to$ \glsterm{dpo}{DPO}) to one (\glsterm{orpo}{ORPO} from base directly). The cost: the
reference-frozen guarantee that \glsterm{dpo}{DPO} inherits from
\cref{eq:dpo-rl-optimum} no longer applies, so \glsterm{orpo}{ORPO} is no longer the
same fixed point as PPO-\glsterm{rlhf}{RLHF} on the same data --- it is a different
optimum that may or may not match.

\paragraph{SimPO --- reference-free, length-normalized.} \glsterm{simpo}{SimPO} drops
the reference and divides the \glsterm{implicit-reward}{implicit reward} by sequence length
\citep[\S 3]{meng2024-simpo}:\footnote{KB excerpt:
\texttt{kb/excerpts/meng2024-simpo.md\#sec-3}.}
\begin{equation}
r_{\mathrm{SimPO}}(x, y)
  \;=\; \frac{\beta}{|y|} \log \pi_\theta(y \mid x)
  \;=\; \frac{\beta}{|y|} \sum_{t=1}^{|y|} \log \pi_\theta(y_t \mid x, y_{<t}).
\label{eq:simpo-reward}
\end{equation}
The full loss adds a target reward margin $\gamma$
\citep[\S 3, the SimPO loss]{meng2024-simpo}:
\begin{equation}
\mathcal{L}_{\mathrm{SimPO}}(\theta)
  = -\, \E_{(x, y_w, y_l) \sim \mathcal{D}}
    \!\left[\log \sigma\!\left(
       \frac{\beta}{|y_w|} \log \pi_\theta(y_w \mid x)
       - \frac{\beta}{|y_l|} \log \pi_\theta(y_l \mid x)
       - \gamma
    \right)\right].
\label{eq:simpo-loss}
\end{equation}
Two changes vs.\ \glsterm{dpo}{DPO}. First, no reference model: the \glsterm{implicit-reward}{implicit reward}
is the length-normalized average log-prob of the response under the
policy alone. Second, a target margin $\gamma > 0$: the chosen
response must beat the rejected one by at least $\gamma$ in implicit
reward, not just by an arbitrary positive amount. \citet[\S 4]{meng2024-simpo}
report \glsterm{simpo}{SimPO} outperforming \glsterm{dpo}{DPO} substantially on length-controlled
benchmarks (AlpacaEval~2, Arena-Hard).\footnote{KB excerpt:
\texttt{kb/excerpts/meng2024-simpo.md\#sec-4-headline}.}

\paragraph{CPO --- policy log-prob with SFT regularizer.} CPO drops the
reference (like \glsterm{simpo}{SimPO} and \glsterm{orpo}{ORPO}) and adds an SFT regularizer on the
chosen response (like \glsterm{orpo}{ORPO}), but keeps the \glsterm{dpo}{DPO} log-sigmoid form
\citep[Xu et al.\ 2024]{rafailov2023-dpo}.
\deepencite{xu2024-cpo §3 — bibkey not yet in references.bib} CPO is
empirically positioned between \glsterm{simpo}{SimPO} and \glsterm{orpo}{ORPO} on the lattice and is
most-cited in machine-translation post-training.

\begin{analogy}
The variant family is best read as one organized axis: \emph{anchored
vs.\ reference-free}. Anchored methods (\glsterm{dpo}{DPO}, \glsterm{ipo}{IPO}, \glsterm{kto}{KTO}) keep
$\pi^{\mathrm{ref}}$ in the loss; the \glsterm{implicit-reward}{implicit reward} is a
\emph{ratio}, and the policy is constrained --- through the algebra of
$\log \pi_\theta / \pi^{\mathrm{ref}}$ --- to live on the support of
the reference. Reference-free methods (\glsterm{orpo}{ORPO}, \glsterm{simpo}{SimPO}, CPO) drop the
reference; the \glsterm{implicit-reward}{implicit reward} becomes an \emph{absolute} quantity
($\log \pi_\theta(y \mid x)$ or its length-normalized version), and
the policy is no longer pinned to the SFT support. Returning to
canonical form: anchored variants inherit the
$\KL(\pi_\theta \| \pi^{\mathrm{ref}})$ regularization of
\cref{eq:dpo-rl-optimum} as a structural property of the implicit-
reward parameterization; reference-free variants do not, and
substitute different mechanisms (\glsterm{simpo}{SimPO}'s length-norm; \glsterm{orpo}{ORPO}'s SFT loss;
CPO's SFT regularizer) to keep the policy on a sensible support
without the reference's constraint. The choice of axis determines what
guarantees you keep --- anchored variants share a fixed point with
PPO-\glsterm{rlhf}{RLHF}, reference-free variants do not --- and what cost you pay ---
anchored variants run a reference forward each step, reference-free
variants do not.
\end{analogy}

\subsection{Length bias and SimPO's normalization fix}
\label{sec:dpo-length}

\glsterm{dpo}{DPO}'s documented \emph{length bias} (longer outputs systematically beat
shorter ones on length-uncontrolled benchmarks) is the most
empirically-robust failure mode of the canonical Eq.~\cref{eq:dpo-loss}
form. The mechanism, articulated by \citet[\S 3]{meng2024-simpo}, is
algebraic: the \glsterm{implicit-reward}{implicit reward}
$\hat r_\theta(x, y) = \beta \log \pi_\theta(y \mid x) /
\pi^{\mathrm{ref}}(y \mid x)$ is a sum over tokens. For preference pairs
with very different lengths, \glsterm{dpo}{DPO} can satisfy the preference by assigning
extreme per-token log-probability shifts on a few tokens of the long
response --- the \glsterm{implicit-reward}{implicit reward} gap grows without the per-token
shifts being individually large. This biases optimization toward the
longer side of any length asymmetry in the preference data.

\glsterm{simpo}{SimPO}'s length-normalized reward \cref{eq:simpo-reward} fixes the
mechanism by construction: dividing by $|y|$ converts the sum-of-tokens
\glsterm{implicit-reward}{implicit reward} into a \emph{mean-of-tokens} \glsterm{implicit-reward}{implicit reward}, which is
length-invariant in expectation. \citet[\S 4.2]{meng2024-simpo} report
the length-normalization ablation as the dominant single contributor to
\glsterm{simpo}{SimPO}'s gain over \glsterm{dpo}{DPO} on length-controlled
AlpacaEval~2.\footnote{KB excerpt:
\texttt{kb/excerpts/meng2024-simpo.md\#sec-4-headline}.}

\begin{intuition}
\glsterm{simpo}{SimPO}'s length-normalization is to \glsterm{dpo}{DPO}'s \glsterm{implicit-reward}{implicit reward} what
\glsterm{perplexity}{perplexity} is to log-likelihood. Raw log-likelihood favors short
sequences (each token contributes a negative term to the total);
\glsterm{perplexity}{perplexity} normalizes by length to compare sequences fairly across
sequence-length regimes. Similarly, \glsterm{dpo}{DPO}'s $\log \pi_\theta(y \mid x)$
favors sequences whose token-level log-probs are easy to inflate (the
loss reads a sum); \glsterm{simpo}{SimPO}'s $\log \pi_\theta(y \mid x) / |y|$ normalizes
for length so the \glsterm{implicit-reward}{implicit reward} is comparable across responses of
unequal length. Returning to the canonical math:
\cref{eq:simpo-reward} is the average of $\log \pi_\theta(y_t \mid x,
y_{<t})$ over the trajectory, scaled by $\beta$, so the
preference-margin gradient \cref{eq:dpo-gradient} acquires a
$1 / |y|$ factor that prevents long-response inflation.
\end{intuition}

\subsection{Frontier-2026 deployment: Llama-3, T\"ulu-3, Gemma 3}
\label{sec:dpo-frontier}

The empirical narrative around \glsterm{dpo}{DPO} at frontier scale has settled into
two camps. The Llama-3 / T\"ulu-3 camp \citep{meta-llama3, tulu3} runs
SFT $\to$ \glsterm{rejection-sampling-sft}{rejection-sampling SFT} $\to$ multi-round iterated \glsterm{dpo}{DPO} and
reports that \glsterm{dpo}{DPO} required less compute than PPO and performed as well
or better at the post-training stage \citep[\S 4]{meta-llama3}. T\"ulu-3
\citep{tulu3} adopts the same SFT $\to$ \glsterm{dpo}{DPO} $\to$ \glsterm{rlvr-3}{RLVR}
(\cref{sec:rlvr-grpo}) stack as its open recipe, and explicitly
reproduces the per-step cost-and-quality crossover of
\cref{tab:dpo-vs-ppo-cost}. The Gemma 3 / DeepMind camp
\citep{gemma3} retains a PPO-style \glsterm{rlhf}{RLHF} stage with \glsterm{warm}{WARM}/\glsterm{warp}{WARP}-augmented
reward ensembles \citep{warm-2024, warp-2024} as the production
preference stage, on the argument that ensembled-RM PPO recovers gains
that the \glsterm{dpo}{DPO}-replaces-PPO move sacrifices on creative or open-ended
generation.

The orthogonal axis in 2026 is the relationship between the offline-
preference stage and what surrounds it. \cref{sec:rlaif-cai} walks the
\glsterm{rlaif-3}{RLAIF} / \glsterm{constitutional-ai}{Constitutional AI} \citep{bai2022-cai} family: AI judgments under
a \glsterm{constitution}{constitution} can replace human labels at the preference-data step,
and the resulting AI-preference triples can feed \glsterm{dpo}{DPO} directly with no
PPO loop in between. The composition --- AI-labeled preferences
$\to$ \glsterm{dpo}{DPO} --- is increasingly the default at open-source scale,
because both halves of the pipeline avoid human labeling and on-policy
RL. \cref{sec:rlvr-grpo}'s \glsterm{rlvr-3}{RLVR} / \glsterm{grpo-2}{GRPO} stack is the verifier-domain
alternative: where a programmatic verifier exists, the learned reward
of \cref{sec:rlhf-rm} is replaced by an indicator
$R_{\mathrm{verifier}}(x, y) \in \{0, 1\}$ and the implicit-reward
parameterization of \cref{eq:dpo-implicit-reward} is replaced by a
group-mean baseline. \glsterm{dpo}{DPO} and \glsterm{rlvr-3}{RLVR} are orthogonal: \glsterm{dpo}{DPO} replaces the
\emph{algorithm} (PPO $\to$ classification on preferences); \glsterm{rlvr-3}{RLVR}
replaces the \emph{reward source} (learned RM $\to$ verifier). Modern
post-training pipelines compose both --- T\"ulu-3 runs \glsterm{dpo}{DPO} for general
preferences and \glsterm{grpo-2}{GRPO} for verifier-domain reasoning --- treating them as
complementary stages.

\subsection{Open contradictions}
\label{sec:dpo-contradictions}

Three contradictions are surfaced here and picked up in \cref{sec:rlhf}'s
\S~14 contradiction concentration.

\begin{contradiction}
\textbf{\glsterm{dpo}{DPO}-vs-PPO at frontier scale.} The Llama-3 / T\"ulu-3 view
\citep{meta-llama3, tulu3} is that \glsterm{dpo}{DPO} is \emph{better} than PPO at
frontier scale, not just cheaper: with iterated on-policy preference
collection, the \glsterm{dpo}{DPO}-class pipeline matches or exceeds PPO-class
pipelines on standard alignment benchmarks at lower compute. The
Gemma 3 / DeepMind view \citep{gemma3, warm-2024, warp-2024} is that
PPO + \glsterm{warm}{WARM}/\glsterm{warp}{WARP}-averaged reward ensembles is still the production
choice for general alignment, with the ensembled RM recovering the
proxy-vs-gold gap that single-RM PPO suffered. The disagreement is
partly recipe-specific (different RMs, different preference datasets,
different KL schedules) and partly task-specific: \glsterm{dpo}{DPO} is more
competitive on instruction-following and structured chat;
PPO-with-ensemble retains advantage on creative / open-ended
generation per anecdotal reports. No controlled head-to-head at
frontier scale has been published, and the data, RM, and KL-schedule
axes are entangled in the existing comparisons. Pick up in \S 14.
\end{contradiction}

\begin{contradiction}
\textbf{\glsterm{dpo}{DPO} does not avoid \glsterm{reward-overoptimization}{reward overoptimization}.} The early \glsterm{dpo}{DPO}
narrative was structural: ``no explicit \glsterm{reward-model}{reward model} means no reward
hacking'' --- if there is no $r_\phi$ to overfit, the proxy-vs-gold
gap that drives PPO's reward-hacking pathology cannot exist. The
2024 ``scaling laws for direct alignment'' result
\deepencite{rosset2024-direct-alignment-scaling arXiv:2406.02900 §3
— bibkey not yet in references.bib} undercuts this read. The implicit
reward $\hat r_\theta = \beta \log \pi_\theta / \pi^{\mathrm{ref}}$ of
\cref{eq:dpo-implicit-reward} is itself a learned proxy --- it is
\emph{trained} on the preference data through the maximum-likelihood
fit of \cref{eq:dpo-loss} --- and the gap between $\hat r_\theta$ and
the gold preference grows with $\KL(\pi_\theta \| \pi^{\mathrm{ref}})$
in the same way that the gap between $r_\phi$ and gold preference
grows in PPO-\glsterm{rlhf}{RLHF}. The mechanism is more subtle than ``\glsterm{dpo}{DPO} has no
\glsterm{reward-model}{reward model}'': the policy log-ratio \emph{is} the \glsterm{reward-model}{reward model}, and
the policy hacks itself. The empirical implication: the
$\KL$-anchor regularization is load-bearing for both PPO-\glsterm{rlhf}{RLHF} and
\glsterm{dpo}{DPO}, and the practical mitigations --- on-policy iteration, KL
scheduling, early stopping --- are necessary for \glsterm{dpo}{DPO} too. Pick up
in \S 14.
\end{contradiction}

\begin{contradiction}
\textbf{Length bias is annotator-bias-or-loss-form.} \glsterm{simpo}{SimPO} frames the
length-normalization fix as a loss-form correction
\citep[\S 3]{meng2024-simpo}: \glsterm{dpo}{DPO}'s implicit-reward sum-over-tokens
biases toward longer outputs by construction, and the fix is
algebraic. The competing view is that the length bias originates in
the \emph{preference data}: human (and AI) annotators systematically
prefer longer, more detailed responses, and the bias transfers
through the loss regardless of normalization. Both are partially
right: \glsterm{simpo}{SimPO}'s $1/|y|$ fix demonstrably reduces length bias on
length-controlled benchmarks \citep[\S 4]{meng2024-simpo}, but
remaining length bias persists in the data after normalization. The
two factors have not been cleanly separated. The empirical signature
that would distinguish them --- a length-controlled human preference
collection followed by both \glsterm{dpo}{DPO} and \glsterm{simpo}{SimPO} training, with held-out
length-controlled evaluation --- has not, to our knowledge, been
published at frontier scale.
\end{contradiction}

\subsection{Forward link}
\label{sec:dpo-forward}

\cref{sec:dpo-derivation} establishes the inversion of
\cref{eq:rlhf-optimum} into the closed-form classification loss
\cref{eq:dpo-loss}; \cref{sec:dpo-protocol} the per-step
infrastructure delta versus PPO; \cref{sec:dpo-variants} the variant
family organized along reference-anchoring and length-normalization;
\cref{sec:dpo-length} \glsterm{simpo}{SimPO}'s length-bias fix; \cref{sec:dpo-frontier}
the frontier-2026 deployment landscape; \cref{sec:dpo-contradictions}
the three live contradictions concentrated in \S 14. Two next stages
follow. \cref{sec:rlaif-cai} replaces the \emph{preference source}
itself: AI-judge-under-a-\glsterm{constitution}{constitution} preferences feed the same \glsterm{dpo}{DPO}
loss directly, with \cref{eq:dpo-loss} unchanged but the dataset
$\mathcal{D}$ now generated by a preference-labeling \glsterm{llm}{LLM} rather than
human annotators. \cref{sec:rlvr-grpo} replaces the \emph{reward
mechanism}: where a programmatic verifier exists, the learned reward
of \cref{sec:rlhf-rm} and the \glsterm{implicit-reward}{implicit reward} of
\cref{eq:dpo-implicit-reward} are both displaced by an indicator
$R_{\mathrm{verifier}} \in \{0, 1\}$, and the optimization moves from
preference classification on log-ratios back to a critic-free RL loop
(\glsterm{grpo-2}{GRPO}) that exploits the verifier's lack of a proxy-vs-gold gap. The
canonical 2026 post-training stack runs all three stages in sequence:
SFT (\cref{sec:sft}) $\to$ preference-anchored \glsterm{dpo}{DPO} or its variants
(\cref{sec:dpo}) $\to$ verifier-domain \glsterm{rlvr-3}{RLVR} with \glsterm{grpo-2}{GRPO}
(\cref{sec:rlvr-grpo}), with adaptation and merging
(\cref{sec:adaptation}) closing the pipeline.
